\subsection{Thiết kế hai luồng vận hành chính}

Phân hệ RAG có hai luồng vận hành chính: Indexing Flow và Query Flow. Hai luồng này không độc lập hoàn toàn mà sử dụng chung kho dữ liệu, metadata, vector index, graph và trạng thái tài liệu.

Indexing Flow là luồng sản xuất tri thức. Luồng này chạy khi quản trị viên upload tài liệu mới, scan thư mục tài liệu hoặc xử lý lại tài liệu lỗi. Nhiệm vụ của Indexing Flow là biến tài liệu thô thành các đơn vị tri thức có thể truy hồi như full document, chunk, embedding, entity, relation và graph.

Query Flow là luồng tiêu thụ tri thức. Luồng này chạy khi người dùng đặt câu hỏi. Nhiệm vụ của Query Flow là chuẩn hóa câu hỏi, truy hồi bằng chứng từ vector store hoặc graph, xây dựng context và gọi LLM để tạo câu trả lời có căn cứ. Tài liệu Flow report cũng mô tả rõ RAG có hai pipeline chính: Indexing Flow biến tài liệu thô thành chunk, vector, entity, relation và graph; Query Flow biến câu hỏi tự nhiên thành truy vấn, tìm bằng chứng, xây dựng ngữ cảnh và yêu cầu LLM sinh câu trả lời.

\subsubsection{Indexing Flow}

Indexing Flow là luồng sản xuất tri thức, bắt đầu từ tài liệu thô, đi qua các bước đọc file, kiểm tra định dạng, trích xuất nội dung, đưa vào hàng đợi, chunking, trích xuất entity/relation, ghi graph/vector và cập nhật trạng thái tài liệu.

\begin{figure}[H]
    \centering
    \includegraphics[width=\textwidth,height=0.82\textheight,keepaspectratio]{\detokenize{content/source/RAG/flows/kmarag_indexing_flow.png}}
    \caption{Luồng Indexing Flow trong KmaRAG}
    \label{fig:rag-indexing-flow}
\end{figure}

\subsubsection{Query Flow}

Query Flow bắt đầu từ request \texttt{/query/stream}, sau đó hệ thống validate dữ liệu đầu vào, chuẩn hóa tham số, tạo query parameter cho LLM, kiểm tra mode truy vấn và điều hướng sang các nhánh local, global, hybrid, mix, naive hoặc bypass.

\begin{figure}[H]
    \centering
    \includegraphics[width=\textwidth,height=0.82\textheight,keepaspectratio]{\detokenize{content/source/RAG/flows/kmarag_query_flow.png}}
    \caption{Luồng truy vấn trong KmaRAG}
    \label{fig:rag-query-flow}
\end{figure}

Trong quá trình truy vấn, hệ thống không chỉ truy hồi và sinh câu trả lời theo một đường thẳng. Một số điểm kiểm soát được đặt trong luồng, ví dụ nếu thiếu thực thể thì hỏi lại, nếu không có bằng chứng thì không trả lời khẳng định, nếu mâu thuẫn nguồn thì cần reranking hoặc thông báo mức độ không chắc chắn.

Điểm quan trọng là Query Flow chỉ hoạt động tốt nếu Indexing Flow tạo ra dữ liệu sạch, có metadata và có chỉ mục đúng. Nếu tài liệu bị parse lỗi, chunk mất ngữ cảnh hoặc metadata thiếu năm áp dụng, hệ thống có thể truy hồi nhầm nguồn và sinh câu trả lời không đáng tin cậy.
